% Source: Weinberg GR, physical PDF pp. 616--624
%         (printed pp. 588--593 and 595--597; boundary pages are partial),
%         plus restored source-supplements/printed-p594.png.
% Coverage: Section 15.11, The Very Early Universe; equations
%           (15.11.1)--(15.11.28), including all unnumbered displays.
% Figures/tables/footnotes: reference notes 149, 156, and 158--172; no figures
%                          or tables.
% Status: source-reviewed and compile-clean.
% Uncertainties: the final line of (15.11.4) omits the exponent -1/2 carried
%                by the preceding line. Also, (15.11.21) prints
%                (4N/5)^(1/2), whereas (15.11.4) and (15.11.20) give its
%                reciprocal coefficient, (5N/4)^(1/2). Both source forms are
%                preserved and marked VERIFY SOURCE below.
% RESTORED SOURCE: source-supplements/printed-p594.png supplies equations
%                  (15.11.25)--(15.11.27) and their surrounding prose.
% MODERNIZATION CHECK: R_{\mathrm{source}}(t) -> a(t). Weinberg's radiation
% constant a_source is written a_{\mathrm{rad}} to avoid collision with the
% modern scale factor. Coordinate and thermodynamic derivatives are compact.

\subsection{The Very Early Universe}\label{sec:15.11}

The thermal history of the universe was traced in Section~\ref{sec:15.6} back
to an era when the temperature was about \(10^{12}\,\mathrm K\). At this early
time, the universe was filled with particles---photons, leptons, and
antileptons---whose interactions are hopefully weak enough to allow this medium
to be treated as a more or less ideal gas. However, if we look back a little
further, into the first \(0.0001\,\mathrm s\) of cosmic history when the
temperature was above \(10^{12}\,\mathrm K\), we encounter theoretical
problems of a difficulty beyond the range of modern statistical mechanics. At
such temperatures, there will be present in thermal equilibrium copious numbers
of strongly interacting particles---mesons, baryons, and antibaryons---with a
mean interparticle distance less than a typical Compton wavelength. These
particles will be in a state of continual mutual interaction, and cannot
reasonably be expected to obey any simple equation of state.

However, the temptation to try to construct some sort of model of the very
early universe is irresistible. There are in fact two extremely different
simple models that have been widely considered in recent years, and that
reflect two divergent views of the nature of the strongly interacting
particles. Although neither model can be taken seriously in detail, the hope is
that one or the other of these models may come close enough to reality to lead
to useful insights about the very early universe.

The first of these two pictures may be called the \emph{elementary-particle
model}. It is supposed that all particles are made up of a small number of
elementary particles---say, photons, leptons, ``quarks,'' and their
antiparticles. It is further supposed here that at very high temperatures the
forces that bind the elementary particles become negligible, just as the
neutron--proton force becomes cosmologically unimportant at temperatures above
the dissociation temperature of deuterium. Let there be \(\mathcal N\)
different kinds of elementary particles, counting spin states and antiparticles
separately, and counting fermions as \(7/8\) of a particle. (See
Eq.~\eqref{eq:15.6.32}. For instance, if we include only the familiar photons,
leptons, and antileptons, plus three kinds of spin-\(\tfrac12\) quarks and
antiquarks, we have \(\mathcal N=26\).) Then for \(kT\) above the mass of
the heaviest elementary particle, the contents of the universe will behave as
if they consisted of \(\mathcal N/2\) different kinds of blackbody radiation,
with pressure, energy density, and specific entropy given by
\begin{equation}
  3p\simeq\rho\simeq\frac12\mathcal N a_{\mathrm{rad}}T^4,
  \tag{15.11.1}\label{eq:15.11.1}
\end{equation}
\begin{equation}
  \sigma
  \simeq
  \frac{\rho+p}{n_BkT}
  \simeq
  \frac{2a_{\mathrm{rad}}T^3}{3n_Bk}\,\mathcal N,
  \tag{15.11.2}\label{eq:15.11.2}
\end{equation}
where \(n_B\) is the number density of baryons minus antibaryons. (The extra
factor \(1/2\) enters here to cancel the factor 2 in the
Stefan--Boltzmann constant arising from the two photon polarization states.)
For an adiabatic expansion \(\sigma\) is constant, and since \(\sigma\) at
present has the value \eqref{eq:15.5.18}, the temperature in the very early
universe is given by
\begin{equation}
  \frac{T}{T_{\gamma0}}
  =
  \left(\frac{2n_B}{\mathcal N n_{B0}}\right)^{1/3}
  =
  \left(\frac2{\mathcal N}\right)^{1/3}
  \frac{a_0}{a}.
  \tag{15.11.3}\label{eq:15.11.3}
\end{equation}
The relation between the energy density \(\rho\) and the time \(t\) here is
the same as \eqref{eq:15.6.44}, and so
\begin{equation}
  \begin{split}
  t
  &\simeq
  \left(\frac{32\pi G\rho}{3}\right)^{-1/2}
  \simeq
  \left(
    \frac{16\pi G\mathcal N a_{\mathrm{rad}}T^4}{3}
  \right)^{-1/2}\\
  &\simeq
  \left(
    \frac{32\pi G a_{\mathrm{rad}}T_{\gamma0}^4}{3}
  \right)
  \left(\frac{\mathcal N}{2}\right)^{1/6}
  \left(\frac a{a_0}\right)^2.
  \end{split}
  \tag{15.11.4}\label{eq:15.11.4}
\end{equation}
% VERIFY SOURCE: the scan omits -1/2 from the first factor in the final line,
% although the preceding line and dimensional analysis require it.

In contrast, in a \emph{composite-particle model}, it is supposed that there
are no true elementary strongly interacting particles but instead that all
such ``hadrons'' must be regarded as composites of one another. We then face a
question of principle, whether thermodynamic calculations can be carried out
including as ``particles'' only the one absolutely stable hadron, the proton,
or also slowly decaying hadrons, such as the neutron and pi-mesons, or perhaps
\emph{all} resonant hadron states, including such rapidly decaying resonances
as the rho-meson and the ``3--3'' \(\pi\)--\(N\) resonance. It is an
attractive conjecture, that if all resonances were included in our
thermodynamic calculations, then to a first approximation, it might not be
necessary to take any further account of the particles' mutual interactions.
(See Section~\ref{sec:11.9}.) If so, then the contents of the early universe
could be treated as consisting of a great many ideal gases, with
\(\mathcal N(m)\,\dd m\) types between mass \(m\) and \(m+\dd m\). But
what is the function \(\mathcal N(m)\)? The greatest possible contrast with
the elementary-particle model is achieved for a distribution that grows as
fast as possible, that is,
\begin{equation}
  \mathcal N(m)
  \longrightarrow
  Am^{-B}\exp\left(\frac{m}{kT_M}\right)
  \qquad\text{for }m\to\infty,
  \tag{15.11.5}\label{eq:15.11.5}
\end{equation}
where \(A\), \(B\), and \(T_M\) are unknown constants. Thermodynamic
quantities will generally involve integrals of \(\mathcal N(m)\,\dd m\),
with weighting factors that behave like \(\exp(-m/kT)\) for
\(m\to\infty\), so these quantities would not converge for a distribution
function that grew faster than \eqref{eq:15.11.5}, and, even for the
distribution function \eqref{eq:15.11.5}, would not converge for
\(T>T_M\). An ideal-gas model with a number of species given by
Eq.~\eqref{eq:15.11.5} is thus characterized by a \emph{maximum temperature}
\(T_M\). The analysis of secondary particle emission in very high-energy
reactions\hyperref[ch15-ref-158]{\textsuperscript{158}} and the recent
Veneziano model of hadron
interactions\hyperref[ch15-ref-159]{\textsuperscript{159}} both independently
suggest a number of hadron species given by Eq.~\eqref{eq:15.11.5}, with
\(B\) of order 2 to 4, and with \(T_M\) of order
\(1.7\times10^{12}\,\mathrm K\). Leaving aside mesons, leptons, and photons
for the moment, the total energy density, pressure, and baryon number density
are given by the usual Fermi distribution as
\begin{equation}
  \rho
  =
  \hbar^{-3}
  \int\mathcal N(m)\,\dd m
  \int
  \left\{
    \bigl[\exp((E-\mu)/kT)+1\bigr]^{-1}
    +\bigl[\exp((E+\mu)/kT)+1\bigr]^{-1}
  \right\}
  E\,\dd^3p,
  \tag{15.11.6}\label{eq:15.11.6}
\end{equation}
\begin{equation}
  \begin{split}
  p
  =
  \frac13\hbar^{-3}
  \int\mathcal N(m)\,\dd m
  \int
  \Bigl\{
    &\bigl[\exp((E-\mu)/kT)+1\bigr]^{-1}\\
    &+\bigl[\exp((E+\mu)/kT)+1\bigr]^{-1}
  \Bigr\}
  E^{-1}p^2\,\dd^3p,
  \end{split}
  \tag{15.11.7}\label{eq:15.11.7}
\end{equation}
\begin{equation}
  n
  =
  \hbar^{-3}
  \int\mathcal N(m)\,\dd m
  \int
  \left\{
    \bigl[\exp((E-\mu)/kT)+1\bigr]^{-1}
    -\bigl[\exp((E+\mu)/kT)+1\bigr]^{-1}
  \right\}
  \dd^3p,
  \tag{15.11.8}\label{eq:15.11.8}
\end{equation}
where \(E\) is the particle energy \((p^2+m^2)^{1/2}\), and \(\mu\) is
the chemical potential associated with baryon number. The dimensionless
entropy per baryon \(\sigma\) is defined by the second law of thermodynamics
as the integral of the perfect differential
\[
  \dd\sigma
  =
  \frac1{kT}
  \left\{
    \dd\left(\frac\rho n\right)
    +p\,\dd\left(\frac1n\right)
  \right\},
\]
and a straightforward integration gives
\begin{equation}
  \sigma=\frac{\rho+p-\mu n}{nkT}.
  \tag{15.11.9}\label{eq:15.11.9}
\end{equation}
In an adiabatic expansion, \(\rho\) and \(p\) drop from presumably infinite
initial values, while \(\sigma\) must stay constant. The only way that
\(\rho\) and \(p\) can approach infinity while \(\sigma\) remains fixed is
for \(\mu\) to become infinite while \(T\) approaches a finite value
\(T_1\) less than \(T_M\). In this limit, the integrals
\eqref{eq:15.11.6}--\eqref{eq:15.11.8} approach the
values\hyperref[ch15-ref-160]{\textsuperscript{160}}
\begin{equation}
  \begin{split}
  \rho
  \longrightarrow{}&
  A'e^{\mu/kT_M}\mu^{5/2-B}kT_1
  \csc\left(\frac{\pi T_1}{T_M}\right)\\
  &\times
  \left\{
    1
    +\frac{\pi kT_1}{\mu}
      \left(B-\frac52\right)
      \cot\left(\frac{\pi T_1}{T_M}\right)
    +\frac{3kT_M}{2\mu}
      \left(B-\frac14\right)
    +\mathcal O\left(\frac1{\mu^2}\right)
  \right\},
  \end{split}
  \tag{15.11.10}\label{eq:15.11.10}
\end{equation}
\begin{equation}
  p
  \longrightarrow
  A'e^{\mu/kT_M}\mu^{5/2-B}kT_1
  \csc\left(\frac{\pi T_1}{T_M}\right)
  \left\{
    \frac{kT_M}{\mu}
    +\mathcal O\left(\frac1{\mu^2}\right)
  \right\},
  \tag{15.11.11}\label{eq:15.11.11}
\end{equation}
\begin{equation}
  \begin{split}
  \mu n
  \longrightarrow{}&
  A'e^{\mu/kT_M}\mu^{5/2-B}kT_1
  \csc\left(\frac{\pi T_1}{T_M}\right)\\
  &\times
  \left\{
    1
    +\frac{\pi kT_1}{\mu}
      \left(B-\frac32\right)
      \cot\left(\frac{\pi T_1}{T_M}\right)
    +\frac{3kT_M}{2\mu}
      \left(B-\frac14\right)
    +\mathcal O\left(\frac1{\mu^2}\right)
  \right\},
  \end{split}
  \tag{15.11.12}\label{eq:15.11.12}
\end{equation}
where
\[
  A'=(kT_M)^{3/2}\hbar^{-3}(8\pi)^{-1/2}A.
\]
The entropy \eqref{eq:15.11.9} then takes the value
\begin{equation}
  \sigma
  =
  \frac{T_M}{T_1}
  -\pi\cot\left(\frac{\pi T_1}{T_M}\right).
  \tag{15.11.13}\label{eq:15.11.13}
\end{equation}
Since \(\sigma\) is very large, the initial temperature \(T_1\) is very
close to the maximum temperature \(T_M\):
\begin{equation}
  \frac{T_1}{T_M}
  \simeq
  1-\frac1\sigma+\mathcal O\left(\frac1{\sigma^2}\right).
  \tag{15.11.14}\label{eq:15.11.14}
\end{equation}
With a finite initial temperature, the energy density and pressure of mesons,
leptons, and photons is negligible in the limit \(t\to0\) in comparison with
the baryonic contributions \eqref{eq:15.11.10} and
\eqref{eq:15.11.11}, justifying the neglect of all particles but baryons in
the above calculations. The baryon number density must vary as \(a^{-3}\),
so Eqs.~\eqref{eq:15.11.12} and \eqref{eq:15.11.10} give for \(a\to0\)
\begin{equation}
  \mu\longrightarrow3kT_M\lvert\ln a\rvert
  \tag{15.11.15}\label{eq:15.11.15}
\end{equation}
and
\begin{equation}
  \rho\longrightarrow\mu n
  \propto a^{-3}\lvert\ln a\rvert.
  \tag{15.11.16}\label{eq:15.11.16}
\end{equation}
The Einstein field equation \eqref{eq:15.1.20} then has a solution (with
\(k\) neglected) of the
form\hyperref[ch15-ref-160]{\textsuperscript{160}}
\begin{equation}
  a\propto t^{2/3}\lvert\ln t\rvert^{1/2},
  \tag{15.11.17}\label{eq:15.11.17}
\end{equation}
in contrast with the behavior \(a\propto t^{1/2}\) expected in an
elementary-particle model.

How can we distinguish among models of the very early universe? As pointed
out in Section~\ref{sec:15.6}, most of the constituents of the universe were
in thermal equilibrium at temperatures above \(10^{12}\,\mathrm K\), so that
the present contents of the universe for the most part depend only on the
entropy per baryon, and perhaps on the ratio of the lepton numbers to the
baryon number, in the hot early universe. In order to learn something about
the behavior of the universe before the temperature dropped to
\(10^{12}\,\mathrm K\), we need to look for fossils, particles that might
have escaped thermal equilibrium before the temperature dropped below
\(10^{12}\,\mathrm K\).

One such possible fossil particle is the quark, the hypothetical fundamental
particle of the strong interactions.
Zeldovich\hyperref[ch15-ref-161]{\textsuperscript{161}} has estimated, on the
basis of what we have here called an elementary-particle model, that if quarks
really can exist as free particles, the density of leftover quarks, which
escaped fusion into nucleons in the early universe, would be about the same as
the observed present density of gold atoms. Efforts to find quarks in nature
have so far failed, so one can conclude either that free quarks do not exist, or
that the thermal history of the very early universe was very different from
\eqref{eq:15.11.4}.

One other, less hypothetical, fossil particle is the graviton. As shown by
Eq.~\eqref{eq:15.10.42}, a graviton in an imperfect fluid with shear viscosity
\(\eta\) will have a mean free
time\hyperref[ch15-ref-156]{\textsuperscript{156}}
\begin{equation}
  \tau_g=(16\pi G\eta)^{-1}.
  \tag{15.11.18}\label{eq:15.11.18}
\end{equation}
If \(\tau_g\) is not much longer than the expansion time \(t\), then the
transport of momentum by gravitons will give the medium a viscosity
\begin{equation}
  \eta=\frac4{15}a_{\mathrm{rad}}T^4\tau_g.
  \tag{15.11.19}\label{eq:15.11.19}
\end{equation}
Eliminating \(\eta\) from these two equations then
gives\hyperref[ch15-ref-149]{\textsuperscript{149}}
\begin{equation}
  \tau_g
  =
  \left(
    \frac{64}{15}\pi G a_{\mathrm{rad}}T^4
  \right)^{-1/2}.
  \tag{15.11.20}\label{eq:15.11.20}
\end{equation}
In an elementary-particle model, Eqs.~\eqref{eq:15.11.20} and
\eqref{eq:15.11.4} give the ratio of the graviton mean free time to the
expansion time as
\begin{equation}
  \frac{\tau_g}{t}
  =
  \left(\frac{4\mathcal N}{5}\right)^{1/2}.
  \tag{15.11.21}\label{eq:15.11.21}
\end{equation}
% VERIFY SOURCE: combining the printed (15.11.4) and (15.11.20) instead gives
% (5\mathcal N/4)^{1/2}; the scan clearly prints (4\mathcal N/5)^{1/2}.
If \(\mathcal N\) is not too large, \(\tau_g\) will be not much greater
than \(t\), so that \eqref{eq:15.11.19} and \eqref{eq:15.11.20} will be
roughly correct, and \(\tau_g\) will thus vanish for \(t\to0\) like \(t\).
In an elementary-particle model, then, the ``optical depth''
\(\int\tau_g^{-1}\,\dd t\) of the very early universe for gravitational
radiation diverges logarithmically for \(t\to0\), and the present universe
should thus contain leftover blackbody gravitational
radiation,\hyperref[ch15-ref-162]{\textsuperscript{162}} with a temperature
\begin{equation}
  T_{g0}
  =
  \frac{(Ta)_{t\to0}}{a_0}
  =
  \left(\frac{\mathcal N}{2}\right)^{-1/3}T_{\gamma0}.
  \tag{15.11.22}\label{eq:15.11.22}
\end{equation}
For instance, if \(\mathcal N=26\) and \(T_{\gamma0}=2.7\,\mathrm K\),
then the present gravitational-radiation background temperature is about
\(0.9\,\mathrm K\). In contrast, in a composite-particle model, the product
\(aT\) vanishes for \(t\to0\), so even if the universe is ``optically
thick'' to gravitational radiation, the present graviton background
temperature would be very much less than the value \eqref{eq:15.11.22}.
Thus the presence or the absence of a gravitational-radiation background with
temperature of order \(1\,\mathrm K\) would provide clear evidence for the
behavior of matter in the very early universe. Unfortunately, there does not
seem to be any way to detect a \(1\,\mathrm K\) gravitational-radiation
background directly.\hyperref[ch15-ref-162]{\textsuperscript{162}} Its most
important effect would be to shorten somewhat the expansion time scale during
the radiation-dominated era, very slightly increasing the cosmic production of
helium.

It may be that the heat represented by the huge entropy per baryon in the
microwave background provides the most useful clue to the very early history of
the universe. Of course, it is possible that this heat was put in at the initial
singularity, in which case we must regard \(\sigma\) as a dimensionless
fundamental constant, like the fine-structure constant. However, it is more
attractive to suppose that the present entropy per baryon was generated by
physical dissipative processes acting in the early, or the very early,
universe.

One such nonadiabatic mechanism for entropy generation is provided by the
phenomenon of \emph{bulk viscosity}. We saw in
Section~\ref{sec:15.10} that the shear viscosity and heat conduction can play
no role in a Robertson--Walker model. The only dissipative effect that can
enter in the energy--momentum tensor for an isotropic homogeneous expansion is
the term in Eq.~\eqref{eq:2.11.21} proportional to the bulk viscosity,
\(\zeta\), which in general coordinates takes the form
\[
  \Delta T^{\mu\nu}
  =
  -\zeta(g^{\mu\nu}+U^\mu U^\nu)\nabla_\lambda U^\lambda,
\]
where \(U^\mu\) is the fluid velocity four-vector. In a
Robertson--Walker model, we have \(\nabla_\lambda U^\lambda=3H\), so the
total energy--momentum tensor here is
\[
  T^{\mu\nu}
  =
  \rho U^\mu U^\nu
  +(p-3\zeta H)(g^{\mu\nu}+U^\mu U^\nu).
\]
The whole effect of the bulk viscosity is thus to replace the pressure \(p\)
with
\begin{equation}
  p^*=p-3\zeta H.
  \tag{15.11.23}\label{eq:15.11.23}
\end{equation}
The bulk viscosity therefore has no effect in the formula,
Eq.~\eqref{eq:15.1.20}, for \(\dot a\) in terms of \(\rho\). However, it
does appear in the energy-conservation equation, which now, in place of
\eqref{eq:15.1.21}, reads
\begin{equation}
  \frac{\dd}{\dd a}(\rho a^3)
  =
  -3p^*a^2
  =
  -3pa^2+9\zeta\dot a\,a.
  \tag{15.11.24}\label{eq:15.11.24}
\end{equation}
Since \(n\propto a^{-3}\), the specific entropy will in general increase at
a rate
\[
  \begin{aligned}
  \dot\sigma
  &\equiv
  \frac1{kT}
  \left[
    \frac{\dd}{\dd t}\left(\frac\rho n\right)
    +p\frac{\dd}{\dd t}\left(\frac1n\right)
  \right]\\
  &=
  \frac{\dot a}{na^2kT}
  \left[
    \frac{\dd}{\dd a}(\rho a^3)
    +p\frac{\dd}{\dd a}(a^3)
  \right],
  \end{aligned}
\]
and Eq.~\eqref{eq:15.11.24} gives this
as\hyperref[ch15-ref-149]{\textsuperscript{149}}
\begin{equation}
  \dot\sigma
  =
  \frac{9\zeta\dot a^2}{nkTa^2}.
  \tag{15.11.25}\label{eq:15.11.25}
\end{equation}
For instance, for a fluid consisting of material particles with a very short
mean free time plus photons with mean free time \(\tau\), the bulk viscosity
is\hyperref[ch15-ref-149]{\textsuperscript{149}}
\begin{equation}
  \zeta
  =
  4a_{\mathrm{rad}}T^4\tau
  \left[
    \frac13-\left(\frac{\partial p}{\partial\rho}\right)_n
  \right]^2,
  \tag{15.11.26}\label{eq:15.11.26}
\end{equation}
and Eq.~\eqref{eq:15.11.25} gives an entropy-production rate
\begin{equation}
  \frac{\dot\sigma}{\sigma}
  =
  \frac{3\tau\dot a^2}{a^2}
  \left[
    1-3\left(\frac{\partial p}{\partial\rho}\right)_n
  \right]^2
  \left(
    \frac{4a_{\mathrm{rad}}T^3}{3nk\sigma}
  \right).
  \tag{15.11.27}\label{eq:15.11.27}
\end{equation}
(For neutrinos, just multiply by a factor of \(7/8\).) The production of
entropy here can be understood as due to the fact that the frequency of a free
photon will vary as \(1/a\) between collisions, whereas the temperature of
the material medium will not vary as \(1/a\) unless
\((\partial p/\partial\rho)_n\) takes the value \(1/3\), so that the
expansion of the universe is continually pulling radiation and matter out of
thermal equilibrium with each
other.\hyperref[ch15-ref-163]{\textsuperscript{163}} However, in an
elementary-particle model \((\partial p/\partial\rho)_n\) will be close to
\(1/3\) in the very early universe, whereas in a composite-particle model,
the photons (or neutrinos) have only a small share of the total entropy, so
that the last factor in \eqref{eq:15.11.27} is small. Estimates of
\(\dot\sigma/\sigma\) do not indicate that bulk viscosity can account for
the high entropy of the present
universe.\hyperref[ch15-ref-149]{\textsuperscript{149}}

If the present entropy of the universe is not due to bulk viscosity, then
perhaps it is produced by the effects of shear viscosity or heat conduction in
an initially anisotropic or inhomogeneous expansion. Indeed, it may be just
these dissipative processes that are responsible for smoothing out initial
anisotropies, and hence producing the high degree of isotropy observed in the
cosmic microwave-radiation background.
Misner\hyperref[ch15-ref-164]{\textsuperscript{164}} has shown that neutrino
viscosity, acting before the temperature dropped to
\(2\times10^{10}\,\mathrm K\), would have reduced the present anisotropy of
the blackbody radiation, produced in an initially homogeneous but anisotropic
expansion, to less than \(0.03\%\). (However, see pp.~525--526.)

Another possible explanation for the observed high entropy per baryon is that
the mean baryon number density really vanishes, as in the theories of
Klein\hyperref[ch15-ref-165]{\textsuperscript{165}} and
Alfv\'en.\hyperref[ch15-ref-166]{\textsuperscript{166}} When (and if) the
temperature was above \(10^{13}\,\mathrm K\), the number density of nucleons
plus antinucleons would have been of order
\[
  n_N+n_{\bar N}
  \sim
  \frac{a_{\mathrm{rad}}T^3}{k}
  \sim
  \sigma n_B
  \sim
  \sigma(n_N-n_{\bar N}),
\]
so if \(\sigma\) has remained constant, the fractional excess of nucleons
over antinucleons in the very early universe would have had the very small
value
\[
  \frac{n_N-n_{\bar N}}{n_N+n_{\bar N}}
  \sim
  \frac1\sigma
  \sim
  10^{-8}\ \text{to}\ 10^{-9}.
\]

In the symmetric cosmology of Klein and Alfv\'en this small nucleon excess is
interpreted as a purely local phenomenon---it is supposed that in other parts
of the universe there was a small antinucleon excess, giving rise at present to
galaxies of antimatter. The detailed calculations of
Omn\`es\hyperref[ch15-ref-166a]{\textsuperscript{166a}} show that reasonable
physical processes in a symmetric plasma of matter and antimatter could have
produced the required small separation of matter and antimatter. Unfortunately,
observational astronomy has not yet provided any definite information as to
whether distant galaxies consist of matter or antimatter. Only the gamma-ray
spectrum provides a hint\hyperref[ch15-ref-166b]{\textsuperscript{166b}} that
antimatter may exist on a cosmological scale.

As long as we have the temerity to speculate about the very early universe, we
may as well carry our speculations back to the very beginning. The Friedmann
solutions, taken literally, indicate that \(a\) vanishes as \(t\to0\), as
\(t^{1/2}\) in the elementary-particle models and as
\(t^{2/3}\lvert\ln t\rvert^{1/2}\) in a composite-particle model. For all we
know, this singularity actually occurs, but it is natural to wonder whether it
can be avoided.

One way to avoid a singularity in the very early universe is for the energy
density \(\rho\) to vanish, owing perhaps to some very short-range attractive
force that overbalances the particles' rest masses. If \(\rho\) vanishes at
some critical value \(a_c\) of the scale factor \(a(t)\), then \(\dot a\)
vanishes at \(a_c\) (or, for finite curvature, near \(a_c\)) so that \(a\)
might have decreased to \(a_c\) before beginning its present increasing
phase.

Even if the energy density is always positive, it is still conceivable that
the universe could escape a general singularity, through anisotropies or
inhomogeneities that invalidate the simple Friedmann solutions.
Penrose\hyperref[ch15-ref-167]{\textsuperscript{167}} and
Hawking\hyperref[ch15-ref-168]{\textsuperscript{168}} have proved a number of
powerful theorems that show that a singularity is inescapable under very
general conditions. For instance, one of Hawking's theorems states that a
singularity is unavoidable, provided that general relativity is valid, that
each point of spacetime has a small neighborhood that no curve with timelike
or null tangents cuts more than once, that the energy--momentum tensor
satisfies the positivity condition
\[
  \left[
    T_{\mu\nu}-\frac12g_{\mu\nu}T^\lambda{}_\lambda
  \right]
  W^\mu W^\nu
  \geq0
\]
for all vectors \(W^\mu\) with \(W^\mu W_\mu<0\), and that there is a point
\(p\) such that all the past-directed timelike geodesics through \(p\) start
converging again within a compact region in the past of \(p\). The last
condition is satisfied if there is enough matter to make the world lines
through \(p\) converge in the past, and Hawking and
Ellis\hyperref[ch15-ref-169]{\textsuperscript{169}} have shown that the cosmic
microwave background does provide enough energy density in the past to satisfy
this condition. However, it is important to note that the Penrose--Hawking
theorems do not say that there is a singularity in the past that involves all
space, as in the Friedmann solutions, but only that there is some singularity
somewhere. The singularity might consist merely of one or more isolated
points, which behave like time-reversed collapsing stars.

Finally, it may be that classical general relativity itself breaks down in the
very early universe. One simple way for this to happen is through the effects
of a cosmological constant, to be discussed in the next chapter. A more
intriguing idea is that quantum effects might become important, invalidating
any purely classical field theory of gravitation. For a system, consisting of
point particles with a mean particle energy \(E\), the relative importance of
gravitational ``radiative corrections'' will be described by a ``gravitational
fine-structure constant''
\[
  \alpha_g=\frac{GE^2}{\hbar c^5},
\]
analogous to the usual electromagnetic fine-structure constant \(1/137\).
Quantum effects will become important when \(\alpha_g\) is of order unity,
or when \(E\) reaches a critical value, given in c.g.s.\ units by
\begin{equation}
  E_c
  =
  \left(\frac{\hbar c^5}{G}\right)^{1/2}
  =
  1.22\times10^{28}\,\mathrm{eV},
  \tag{15.11.28}\label{eq:15.11.28}
\end{equation}
corresponding to a temperature \(1.4\times10^{32}\,\mathrm K\). The
temperature in a composite-particle model never gets anywhere near so high,
but \(kT\) would have been greater than \(E_c\) at the very beginning of a
Friedmann universe composed of a finite number \(\mathcal N\) of species of
elementary particles. In fact, a great many other quantum effects become
important at that time.\hyperref[ch15-ref-170]{\textsuperscript{170}} For
instance, the oscillation rate of a typical particle wave function at
temperature \(T\) is \(kT/\hbar\), while the expansion rate of the universe
at this time is given by Eq.~\eqref{eq:15.11.4} as
\[
  H
  =
  \frac1{2t}
  =
  \left(
    \frac{4\pi G\mathcal N a_{\mathrm{rad}}T^4}{3}
  \right)^{1/2}.
\]
Recalling that
\[
  a_{\mathrm{rad}}=\frac{\pi^2k^4}{15\hbar^3},
\]
the ratio of these rates is
\[
  \frac{kT/\hbar}{H}
  =
  \left(
    \frac{45\hbar}{4\pi^3\mathcal N G(kT)^2}
  \right)^{1/2}
  =
  \left(\frac{45}{4\pi^3\mathcal N}\right)^{1/2}
  \left(\frac{E_c}{kT}\right).
\]
Hence for temperatures above the critical temperature \(E_c/k\), the wave
functions of typical particles have an oscillation rate slower than the
expansion rate of the universe, so that no classical or semiclassical
description can be applied to the particles in thermal equilibrium at that
time.

The consideration of first things naturally leads to speculation about last
things.\hyperref[ch15-ref-171]{\textsuperscript{171}} We saw in
Section~\ref{sec:15.1} that a Friedmann universe with \(k=+1\) will
eventually cease its expansion and begin to contract. Taken literally, such
models require that a singularity with \(a=0\) will be reached at some finite
time in the future, of order \(75\times10^9\) years for
\(H_0=75\,\mathrm{km\,s^{-1}\,Mpc^{-1}}\) and \(q_0=1\). However, if it is
possible to escape a general singularity in the past, either through negative
energy densities, anisotropies, or quantum effects, then presumably it ought
to be possible to escape the general singularity in our future. In this case,
we might suppose that the universe undergoes an oscillation, with periods of
contraction and expansion succeeding one another eternally.

Could this oscillation be periodic? That is, can we restore a steady-state
picture of the universe, by viewing cosmic history on a sufficiently grand time
scale? One obvious objection is that entropy is presumably created, not
destroyed, in each cycle. It has been suggested that entropy might be destroyed
during contracting phases,\hyperref[ch15-ref-172]{\textsuperscript{172}}
because it is the expansion of the universe that, by providing a heat sink,
sets the direction of time's arrow in thermodynamic processes. However, there
is no detailed model that describes how this can come about. In particular, it
is hard to see how time's arrow could be reversed just at the moment when
\(a(t)\) reaches its maximum value, at which time the background-radiation
temperature is so low, of order \(1\,\mathrm K\), that it can hardly affect
terrestrial processes. If somehow or other the second law of thermodynamics
can be evaded, and the universe really does expand and contract periodically,
then any particles that are not brought into thermal equilibrium during the
contracted phase, such (perhaps) as gravitons or neutrinos, would have to be
present in large numbers: If \(N\) particles are produced in a given comoving
volume during each cycle, and the probability that one of these particles is
absorbed in one cycle is \(P\), then there would have to be a mean number
\(N/P\) of particles in this volume in order to maintain a more or less
constant population. It is therefore not out of the question that some day we
may detect remnants of previous cycles of the history of the universe. For the
present, however, such matters remain at the furthest bounds of cosmological
speculation.
